\section{A nodal divisor and its exponent diagram}
\label{sec:node}

We now bound the initial exponents of sections of an ample line bundle by
intersecting their divisors with one fixed nodal curve.  Varying the weight
in the two branch coordinates produces a family of quadrilaterals.  The
next subsection will compress these diagrams to triangles.
The variable Newton--Okounkov quadrilaterals of Ciliberto, Farnik,
K{\"u}ronya, Lozovanu, Ro\'e, and Shramov
\cite[Section~5.1 and Corollary~5.8]{CFKLRS2016} suggested this
weighted-diagram approach.  The nodal bound on an arbitrary polarized
surface and the diagonal compression used here are proved directly;
the plane corollary in that source is not a premise.

\begin{lemma}[A nodal ample divisor]
\label{lem:node}
Let $S$ be a smooth integral complex projective surface and let $L$ be an
ample line bundle.  There exist a positive integer $d$, an integral divisor
$V\in|dL|$, and a point $p\in V$ such that $V$ is smooth away from $p$ and
has two smooth transverse analytic branches at $p$.
\end{lemma}

\begin{proof}
Choose an integer $e>0$ such that $eL$ is very ample, embed $S$ in
$\PP^N$ using this line bundle, and choose any $p\in S$.  Take homogeneous
coordinates with $p=[1:0:\cdots:0]$.  The restrictions of $X_1,\ldots,X_N$
generate $\mathcal I_p\otimes\mathcal O_S(eL)$: away from $p$ at least one
is nonzero, and near $p$ their ratios with $X_0$ generate the maximal ideal
because the embedding is a closed immersion.

Let $\pi:X=\operatorname{Bl}_pS\to S$ be the blow-up, with exceptional
curve $E$.  The generators just described give a surjection onto the
twisted Rees algebra
\[
 \operatorname{Sym}^{\bullet}(\mathcal O_S^{\oplus N})
 \longrightarrow
 \bigoplus_{n\ge0}\mathcal I_p^n\otimes\mathcal O_S(neL).
\]
Taking relative $\operatorname{Proj}$ gives a closed immersion
$X\hookrightarrow S\times\PP^{N-1}$, the closure of the graph of projection from $p$.
The second factor pulls $\mathcal O(1)$ back to
$\mathcal O_X(e\pi^*L-E)$; see also \cite[Tag 01OF]{Stacks}.
After embedding the first factor in $\PP^N$, the very ample bundle
$\mathcal O(1,1)$ therefore restricts to
\[
 \mathcal A=\mathcal O_X(2e\pi^*L-E).
\]
In particular, $\mathcal A^{\otimes2}=\mathcal O_X(4e\pi^*L-2E)$ is
very ample.

The point blow-up $X$ is a smooth integral projective surface.  By
Bertini, a general divisor $\widetilde V$ in this last linear system
is smooth and integral \cite[Tags 0FD6 and 0G4F]{Stacks}.  We may also
require that it meet $E$ transversely at two distinct points.  Indeed,
$\mathcal A|_E\simeq\mathcal O_{\PP^1}(1)$, and the restrictions of its
global sections generate this bundle, hence span its full two-dimensional
space of sections.  Products of these sections show that
\[
 H^0(X,\mathcal A^{\otimes2})
 \longrightarrow H^0(E,\mathcal O_{\PP^1}(2))
\]
is surjective.  Having two distinct zeros on $E$ is consequently a
nonempty open condition, which can be imposed along with the Bertini
conditions.

Set $d=4e$ and $V=\pi(\widetilde V)$.  Pushing down the divisor class
shows that $V\in|dL|$, and $V$ is integral because $\widetilde V$ is.
The blow-down is an isomorphism away from $E$, so $V$ is smooth away
from $p$.  In a local blow-up chart the map has the form
$(a_0,u_0)\mapsto(a_0,a_0u_0)$ with $E=\{a_0=0\}$.
A curve transverse to $E$ has $a_0$ as a local parameter, and its image
is a smooth branch.  The two different points of $\widetilde V\cap E$
give different tangent directions.  Thus the two image branches are
smooth and transverse at $p$.
\end{proof}

Fix $d,V,p$ as in Lemma~\ref{lem:node}, and set $H=L^2$.
Choose holomorphic coordinates $(x,z)$ at $p$ in which the two branches
are the coordinate axes, and choose a nonvanishing local frame of $L$.
In this frame a defining section $\tau\in H^0(S,dL)$ of $V$ has the form
\begin{equation}
 \tau=u(x,z)xz,\qquad u(0,0)\ne0.
 \label{eq:node-unit}
\end{equation}
Here and below a section is identified with its local function in the
chosen frame and its tensor powers.

Choose once and for all $T>1$ so large that
\begin{equation}
 (1+t)^2>d^2Ht\qquad\text{for every }t>T.
 \label{eq:weight-range}
\end{equation}
Such a choice is possible because the left side is quadratic in $t$.
For $t>T$, put
\begin{equation}
 \begin{gathered}
 c=\frac1d,\qquad W=dH,\qquad
 a=\frac{Wt}{1+t},\qquad b=\frac{W}{1+t},\\
 Q_t=\conv\{(0,0),(a,0),(c,c),(0,b)\}.
 \end{gathered}
 \label{eq:quadrilateral}
\end{equation}
Recall that $\nu_t(f)$ denotes the first nonzero Taylor exponent of $f$
when exponents $(\alpha,\beta)$ are ordered by
$(\alpha+t\beta,\beta)$.  The inequality \eqref{eq:weight-range} is
equivalent to $c/a+c/b>1$.  Thus $(c,c)$ lies beyond the side joining
$(a,0)$ and $(0,b)$, and the four vertices in
\eqref{eq:quadrilateral} occur in the displayed cyclic order.

\begin{proposition}[The quadrilateral bound]
\label{prop:quadrilateral}
For the surface, line bundle, nodal divisor, coordinates and frame fixed
above, every integer $k\ge1$, every $t>T$, and every nonzero
$s\in H^0(S,kL)$ satisfy
\[
 \nu_t(s)\in kQ_t.
\]
Moreover, $H^0(S,kL)$ has a basis with pairwise distinct initial exponents
in $kQ_t$; the basis is empty if this section space is zero.
\end{proposition}

\begin{proof}
First suppose that the divisor $D$ of $s$ does not contain $V$, and write
$f(x,z)$ for its local function.  Neither $f(x,0)$ nor $f(0,z)$ vanishes
identically.  To see this, $D\cap V$ is a proper closed subset of the
integral algebraic curve $V$, hence is finite.  An identically vanishing
branch restriction would put an entire analytic branch in this finite
set.  This argument uses the global integrality of $V$ without asserting
that its nodal analytic germ is irreducible.

Let $\ell_x=\ord_0f(x,0)$ and $\ell_z=\ord_0f(0,z)$, and let $v$ be
the weight of $\nu_t(f)$.  The pure terms defining these two finite
orders have weights $\ell_x$ and $t\ell_z$, respectively.  Therefore
$\ell_x\ge v$ and $\ell_z\ge v/t$.  Intersection multiplicity is additive
over the two smooth branches, so
\begin{equation}
 v\left(1+\frac1t\right)
 \le \ell_x+\ell_z
 = I_p(D,V)
 \le D\cdot V=kdH.
 \label{eq:branch-bound}
\end{equation}
The analytic computation gives the algebraic local intersection number:
both are the colength of $(f,xz)$ in the common completed local ring
$\C[[x,z]]$.  The final inequality in \eqref{eq:branch-bound} follows
from nonnegativity of all the other local intersection numbers.
Consequently $v\le ka$, and
\begin{equation}
 \nu_t(s)\in k\Delta_t,\qquad
 \Delta_t=\conv\{(0,0),(a,0),(0,b)\},
 \label{eq:residual-triangle}
\end{equation}
since $a=tb$.

For a general nonzero $s$, let $j$ be the coefficient of $V$ in its
divisor.  Dividing by the defining section of this effective Cartier
divisor gives
\[
 s=\tau^j\rho,\qquad
 \rho\in H^0(S,nL),\qquad n=k-jd,
\]
where the effective divisor $R$ of $\rho$ does not contain $V$.
Intersecting with the ample line bundle gives
$nH=L\cdot R\ge0$, so $n\ge0$.  If $n=0$, then $R=0$ by ampleness
and $\rho$ is a nonzero constant, since $S$ is integral and projective.
If $n>0$, the preceding argument gives $\nu_t(\rho)\in n\Delta_t$.

The order on exponents is compatible with addition, so initial exponents
add under multiplication.  A unit has initial exponent $(0,0)$.
Equation~\eqref{eq:node-unit} therefore yields
\[
 \nu_t(s)=(j,j)+\nu_t(\rho).
\]
When $n>0$, this implies
\[
 \frac{\nu_t(s)}{k}
 =\frac{n}{k}\,\frac{\nu_t(\rho)}{n}
  +\frac{jd}{k}(c,c)\in Q_t,
\]
because the two coefficients are nonnegative and sum to one.  When
$n=0$, the normalized exponent equals $(c,c)$ directly.

Finally, restriction to the analytic chart is injective on global
sections, by the identity theorem.  Apply the finite coefficient
elimination in Proposition~\ref{prop:initial-transfer} to this
finite-dimensional space of germs.  It gives a basis with distinct
initial exponents, each of which satisfies the bound just proved.
\end{proof}

Proposition~\ref{prop:quadrilateral} controls the initial exponents; it
makes no claim about the remaining Taylor coefficients.  We next use
slices parallel to $(1,1)$ to replace its quadrilateral by a triangle
for the purpose of testing jets.
